\section*{第 \ref{ch:b3:structural-analysis} 章——结构分析}

\begin{solution}{exo:b3:structural-analysis:1}
极的倾伏向为 $40 + 180 = \ang{220}$，倾伏角为 $90 - 30 = \ang{60}$。
该面倾角为 $90 - 20 = \ang{70}$，倾向为 $315 - 180 = \ang{135}$：$135/70$。
\end{solution}

\begin{solution}{exo:b3:structural-analysis:2}
$\theta = \ang{60}$。等角投影：$10\tan\ang{30} = \qty{5.77}{cm}$；
等积投影：$10\sqrt2\sin\ang{30} = \qty{7.07}{cm}$。
\end{solution}

\begin{solution}{exo:b3:structural-analysis:3}
$\vec n_1 = (\sin\ang{30}, 0, \cos\ang{30})$, $\vec n_2 = (-\sin\ang{30}, 0,
\cos\ang{30})$，$\vec n_1\cdot\vec n_2 = -0.25 + 0.75 = 0.5$：两极相距 
\ang{60}，翼间角为 \ang{120}，属开阔褶皱。
\end{solution}

\begin{solution}{exo:b3:structural-analysis:4}
$\vec d = (\cos\ang{60}, 0, -\sin\ang{60}) = (0.5, 0, -0.866)$ 和 $\vec u =
0.707\,\vec s + 0.707\,\vec d = (0.354, 0.707, -0.612)$：倾伏向
$\operatorname{atan2}(0.354, 0.707) = \ang{26.6}$，倾伏角 $\arcsin0.612 =
\ang{37.8}$。侧伏角介于 \ang{-90}（正断）与 0（左行）之间：一个
斜向正断层，兼具左行分量。
\end{solution}

\begin{solution}{exo:b3:structural-analysis:5}
\begin{gather*}
  \vec n_1 = (0.5567, 0.3214, 0.7660),\qquad \vec n_2 = (-0.4698, -0.1710, 0.8660),\\
  \vec n_1\times\vec n_2 = (0.4093, -0.8420, 0.0558),
\end{gather*}
长度为 $0.9379$。其相反向量 $-\vec n_1\times\vec n_2$ 朝下：
倾伏向 \ang{334.1}，倾伏角 $\arcsin(0.0558/0.9379) = \ang{3.4}$。$\vec n_1\cdot\vec
n_2 = 0.3469$，法线间夹角 \ang{69.7}，翼间角 \ang{110.3}：
属开阔褶皱。
\end{solution}

\begin{solution}{exo:b3:structural-analysis:6}
倾伏角大于 \ang{60} 意味着 $\theta < \ang{30}$。等积投影：$r^2 =
2\sin^2\ang{15} = 0.134$，占网图的 \qty{13.4}{\percent}；等角投影：
$\tan^2\ang{15} = 0.072$，\qty{7.2}{\percent}。在球面上，该球冠占
半球的 $1 - \cos\ang{30} = 0.134$：只有等积网
能如实显示其真实份额，因此极点在其上计数。
\end{solution}

\begin{solution}{exo:b3:structural-analysis:7}
\begin{gather*}
  T_{xx} = \frac{2\sin^2\ang{20}}3 = 0.078,\qquad
  T_{zz} = \frac{2\cos^2\ang{20} + 1}3 = 0.922,\\
  T_{xz} = -\frac{\sin\ang{20}\cos\ang{20} - \sin\ang{20}\cos\ang{20}}3 = 0,
\end{gather*}
所有含 $y$ 的项均为 0。本征值为 $0.922$（竖直方向），
$0.078$（东）与 $0$（北）：\omterm{def:b3:structural-analysis:pi}{β 轴}水平，呈南北向。
三个极点位于同一平面内，故都没有沿其法线方向的分量，
且 $\vec u^{\mathsf T}T\vec u = 0$ 对沿其法线的 $\vec u$ 成立。
\end{solution}

\begin{solution}{exo:b3:structural-analysis:8}
$a = 2$，$b = 1.5$，$k = 1/0.5 = 2$：收缩。$Z = 1/(1.8\times1.1) =
0.505$，$a = 1.636$，$b = 2.178$，$k = 0.636/1.178 = 0.54$：压扁。
\end{solution}

\begin{solution}{exo:b3:structural-analysis:9}
$S_1/S_2 = 2.25$ 且 $S_1S_2 = 1$：$S_1 = 1.5$，$S_2 = 0.667$；长轴
伸长了 \qty{50}{\percent}。孔洞两半轴分别为 $1.5\times0.4 =
\qty{0.60}{mm}$ 与 $0.667\times0.4 = \qty{0.27}{mm}$。
\end{solution}

\begin{solution}{exo:b3:structural-analysis:10}
等角投影：沿径向 $1/(2\cos^2\ang{30}) = 0.667$，且
沿周向 $\tan\ang{30}/\sin\ang{60} = 0.577/0.866 = 0.667$：两者相等，即保角。
等积投影：$\mathrm dr/(R\,\mathrm d\theta) = (\sqrt2/2)\cos\ang{30} = 0.612$
与 $\sqrt2\sin\ang{30}/\sin\ang{60} = 0.707/0.866 = 0.816$；二者之积为
$0.500$，即定理中的面积因子。小圆变为两轴之比
为 $0.816/0.612 = 1.33$ 的椭圆，沿周向的轴较长。
\end{solution}

\begin{solution}{exo:b3:structural-analysis:11}
断层直立，故其法线水平，且垂直于
竖直的 $\sigma_2$；由定理 (b) 部分，滑动
垂直于 $\sigma_2$，即水平方向，对每个 $\Phi$ 均成立。取 $\sigma_1 = 1$ 沿北，
$\sigma_3 = 0$ 沿东：$S\vec n = (0, -0.342, 0)$，$\vec n\cdot S\vec n =
0.117$，$\vec\tau = (0, -0.342, 0) - 0.117\,(0.940, -0.342, 0) = (-0.110,
-0.302, 0)$，$-\vec\tau = (0.110, 0.302, 0)$ 指向 \ang{020}。
位于 $\vec n$ 所指一侧（东侧）的断块沿
走向向北北东运动：从西侧看，它向左运动，为左行断层（侧伏角
为 0）。$\sigma_2$ 未参与计算。
\end{solution}

\begin{solution}{exo:b3:structural-analysis:12}
$S = (7.5 - 6.0)/1.0 = \qty{1.5}{km}$；$h = 6.0/1.5 = \qty{4.0}{km}$，位于
下层位之下。剖面原长为 $20 + 1.5 = \qty{21.5}{km}$，
长度变化 $-1.5/21.5 = \qty{-7.0}{\percent}$。
\end{solution}

\begin{solution}{pb:b3:structural-analysis:1}
\textbf{1.} 极的倾伏向为 \ang{255}，倾伏角为 \ang{77}；$\theta =
\ang{13}$，$r = 10\sqrt2\sin\ang{6.5} = \qty{1.60}{cm}$，距中心
朝 \ang{255} 方向。

\textbf{2.} $\vec n = (\sin\ang{13}\sin\ang{75}, \sin\ang{13}\cos\ang{75},
\cos\ang{13}) = (0.2173, 0.0582, 0.9744)$。六个近似平行的单位矢量
相加，其和接近 6：$|\vec S_E| = 5.9974$。

\textbf{3.} 东翼：$\bar\alpha = \operatorname{atan2}(1.2583, 0.2017) =
\ang{80.9}$，$\bar\delta = \arccos(5.8604/5.9974) = \ang{12.3}$；$\bar R
= 5.9974/6 = 0.99956$，离散度为 $\sqrt{2\times0.00044} =
\qty{0.0297}{rad}$，即 \ang{1.7}。西翼：$\bar\alpha =
\operatorname{atan2}(-1.2445, 0.1886) = \ang{278.6}$，$\bar\delta =
\arccos(5.8636/5.9972) = \ang{12.1}$；$\bar R = 0.99953$，离散度 \ang{1.8}。

\textbf{4.}
\[
  \vec S_E\times\vec S_W = (0.0774, -14.6714, 0.4883),
\]
其模为 $14.680$。其反向朝下：倾伏向为 $\operatorname{atan2}(-0.0774, 14.6714) = \ang{359.7}$，且
倾伏角为 $\arcsin(0.4883/14.680) = \ang{1.9}$。

\textbf{5.} 两平均法向量夹角的余弦为
$32.835/35.968 = 0.9129$，即两和向量的点积除以其
模长：两法向量夹角为 \ang{24.1}，翼间角为 \ang{155.9}：
属开阔褶皱，接近平缓。

\textbf{6.}
\begin{gather*}
  \vec m_E = (0.2098, 0.0336, 0.9772),\qquad \vec m_W = (-0.2075, 0.0314, 0.9777),\\
  \vec m_E - \vec m_W = (0.4173, 0.0022, -0.0005):
\end{gather*}
水平法向量朝向 \ang{089.7}。轴面竖直，走向
\ang{359.7}，实际为南北向：褶皱直立。

\textbf{7.} 环带即极点为 \omterm{def:b3:structural-analysis:pi}{β 轴}的平面：倾向
为 \ang{179.7}，倾角 $90 - 1.9 = \ang{88.1}$，几乎为东西向的竖直
平面。每一翼都含褶皱轴，故每个法向量均与之垂直，
并位于与之垂直的平面内。

\textbf{8.} $0.04400 + 0.00144 + 0.95456 = 1.00000$。$\vec b^{\mathsf T}T\vec
b = 0.00038$，是最小本征值：$\vec b$ 在舍入误差内即
特征向量 $\vec e_3$，而过两个平均极点的环带即
为全部十二个读数的最佳拟合环带。

\textbf{9.} $\arcsin\sqrt{0.00038} = \arcsin0.0195 = \ang{1.1}$。

\textbf{10.} 沿环带，各极点距竖直方向约 \ang{12.2}，
分居两侧；$\sin^2\ang{12.2} = 0.045$ 为其
东西向分量的均方值，即 $\lambda_2$。

\textbf{11.} 是：极点平均偏离环带 \ang{1.1}，
不超过单翼读数绕其均值的离散度 \ang{1.7}。

\textbf{12.} $\ang{1.1}/\sqrt{12} = \ang{0.3}$：倾伏角为 \ang{1.9} $\pm$
\ang{0.3}。

\textbf{13.} 侧伏角为 \ang{-90}：纯正断层，属倾滑。按安德森理论：
$\sigma_1$ 竖直，$\sigma_3$ 沿倾向水平，
\ang{120}--\ang{300}，$\sigma_2$ 沿走向水平，
\ang{030}--\ang{210}。

\textbf{14.} $120/60$ 的法向量位于走向为
\ang{120} 的竖直面内，该面包含 $\vec e_1$ 与 $\vec e_3$：$n_2 = 0$。由 (b) 小题可知，
滑动垂直于 $\sigma_2$（即走向），对所有 $\Phi$ 均成立：属
倾滑。该断层不能说明 $\Phi$ 的任何信息；公式中 $n_2$ 与
$\vec e_2\cdot\vec w$ 同时为零，故分子与分母皆为零。

\textbf{15.} $\vec n = (\sin\ang{55}\sin\ang{150}, \sin\ang{55}\cos\ang{150},
\cos\ang{55}) = (0.4096, -0.7094, 0.5736)$; $n_1 = 0.574$; $n_2 = 0.5\times
0.4096 + 0.866\times(-0.7094) = -0.410$。

\textbf{16.} $\vec e_1\cdot\vec w = 0.198$，$\vec e_2\cdot\vec w = 0.5\times
0.771 + 0.866\times0.605 = 0.909$，故 $\Phi = -(0.574\times0.198)/(-0.410\times
0.909) = 0.1137/0.3727 = 0.305$。

\textbf{17.} $(0.305 + 0.294 + 0.301 + 0.306)/4 = 0.30$。侧伏角
已取整至度；在这些断层上侧伏角每变化一度，$\Phi$ 约变化 0.02，
故半度的取整误差足以解释这一离散。

\textbf{18.} $\sigma_2 - \sigma_3 = 0.3(\sigma_1 - \sigma_3)$：$\sigma_2$
更接近 $\sigma_3$ 而非 $\sigma_1$。两个水平应力均
远小于竖直应力；该区域主要向 \ang{120} 方向伸展，
沿 \ang{030} 方向则略有伸展。

\textbf{19.} 取 $\vec n = (0.750, -0.433, 0.500)$，滑动 $\vec u =
(0.433, -0.250, -0.866)$，则 $(\vec n - \vec u)/\sqrt2$ 倾伏角为 \ang{75}，倾伏向
为 \ang{300}，记为 P；$(\vec n + \vec u)/\sqrt2$ 倾伏角为 \ang{15}，倾伏向为 \ang{120}，
记为 T。二者与 $\sigma_1$ 和 $\sigma_3$ 同处一竖直面，但
各偏离它们 \ang{15}，因为断层与
$\sigma_1$ 的夹角为 \ang{30}，而非 \ang{45}。

\textbf{20.} $\tan a = \tan\ang{12.3}\times\cos(\ang{90} - \ang{80.9}) =
0.2180\times0.9874 = 0.2153$：$a = \ang{12.1}$。

\textbf{21.} $L_0 = 8000/\cos\ang{12} = \qty{8178.7}{m}$，$S =
\qty{178.7}{m}$，相对变化量为 $\cos\ang{12} - 1 =
\qty{-2.2}{\percent}$。

\textbf{22.} $4000\tan\ang{12} = \qty{850.2}{m}$；$A = \tfrac12\times8000\times
850.2 = \qty{3.40e6}{m^2}$。

\textbf{23.} $h = \num{3.40e6}/178.7 = \qty{19.0}{km}$。若在覆盖层底部存在滑脱面，
则 $h$ 仅为数百米，多余面积
约小五十倍：这些褶皱在该处并不滑脱。两种可能是
基底与盖层一同褶皱，或向斜沉降至原始水平之下，
这使 $A$ 变小；仅凭剖面无法判定。

\textbf{24.} 水平断距 $200/\tan\ang{60} = \qty{115.5}{m}$，滑距 $200/\sin\ang{60}
= \qty{230.9}{m}$。

\textbf{25.} $\tan a = \tan\ang{60}\cos\ang{30} = 1.500$：$a = \ang{56.3}$，且
视水平断距为 $200/1.5 = \qty{133.3}{m}$。断层的滑动
沿倾向朝下，即朝 \ang{120} 方向：沿 \ang{120} 的剖面显示真实
长度，应为待恢复的剖面。

\textbf{26.} 断层切割褶皱，故其较新，应先恢复，
将上盘沿断层面向上回移 \qty{230.9}{m}，随后
将褶皱展开。\omterm{def:b3:structural-analysis:restoration}{固定线}应置于向斜的轴面上，
位于下盘，那里对称褶皱的各层之间未发生滑移。

\textbf{27.} 瓦尔登褶皱轴的倾伏角为
\textbf{\ang[text-series-to-math=true]{1.9}}，向北倾伏（倾伏向
\ang{359.7}）：区域内设定的 \ang{2}，恢复后
为十二个读数所允许的 \ang{0.3}。
\end{solution}
